\documentclass[12pt,a4paper]{report}
\usepackage[utf8]{inputenc}
\usepackage[T1]{fontenc}
\usepackage[french]{babel}
\usepackage{geometry}
\usepackage{graphicx}
\usepackage{float}
\usepackage{amsmath}
\usepackage{hyperref}
\usepackage{booktabs}
\usepackage{array}
\usepackage{colortbl}
\usepackage{xcolor}
\usepackage{listings}
\usepackage{enumitem}
\usepackage{titlesec}
\usepackage{fancyhdr}
\usepackage{multirow}
\usepackage{longtable}
\usepackage{tikz}

\usepackage{pgf-umlcd}
\usepackage{caption}
\usepackage{subcaption}
\usepackage{adjustbox}

% TikZ libraries
\usetikzlibrary{arrows.meta, shapes, positioning, calc, decorations.pathreplacing, shadows.blur, backgrounds}

% Configuration des marges
\geometry{left=2.5cm,right=2.5cm,top=2.5cm,bottom=2.5cm}

% Couleurs professionnelles
\definecolor{primaryBlue}{RGB}{0,84,147}
\definecolor{secondaryBlue}{RGB}{0,112,192}
\definecolor{accentGreen}{RGB}{0,128,0}
\definecolor{accentOrange}{RGB}{255,140,0}
\definecolor{lightGray}{RGB}{240,240,240}
\definecolor{darkGray}{RGB}{80,80,80}
\definecolor{tableHeader}{RGB}{0,84,147}
\definecolor{codeBg}{RGB}{250,250,250}

% Style des titres
\titleformat{\chapter}{\normalfont\huge\bfseries\color{primaryBlue}}{\thechapter}{1em}{}
\titleformat{\section}{\normalfont\Large\bfseries\color{secondaryBlue}}{\thesection}{1em}{}
\titleformat{\subsection}{\normalfont\large\bfseries\color{darkGray}}{\thesubsection}{1em}{}

% Style des listings
\lstset{
    basicstyle=\ttfamily\small,
    keywordstyle=\color{primaryBlue}\bfseries,
    commentstyle=\color{accentGreen},
    stringstyle=\color{accentOrange},
    numbers=left,
    numberstyle=\tiny\color{gray},
    frame=single,
    breaklines=true,
    backgroundcolor=\color{codeBg},
    rulecolor=\color{lightGray},
    tabsize=2,
    showspaces=false,
    showstringspaces=false
}

% En-tête et pied de page
\pagestyle{fancy}
\fancyhead[L]{\nouppercase{\leftmark}}
\fancyhead[R]{\textcolor{primaryBlue}{RTSP Multi-Channel System}}
\fancyfoot[C]{\thepage}
\fancyfoot[R]{\textcolor{gray}{Ayoub Smayen}}
\renewcommand{\headrulewidth}{0.5pt}
\renewcommand{\headrule}{\color{primaryBlue}\hrule width\headwidth height\headrulewidth}

% Hyperliens
\hypersetup{
    colorlinks=true,
    linkcolor=primaryBlue,
    urlcolor=secondaryBlue,
    citecolor=accentGreen
}

% Commandes personnalisées
\newcommand{\code}[1]{\texttt{#1}}
\newcommand{\important}[1]{\textcolor{primaryBlue}{\textbf{#1}}}
\newcommand{\success}[1]{\textcolor{accentGreen}{\textbf{#1}}}
\newcommand{\warning}[1]{\textcolor{accentOrange}{\textbf{#1}}}

% Environnement pour les encadrés
\newenvironment{infobox}[1]{%
    \begin{tcolorbox}[colback=lightGray!20, colframe=primaryBlue, title=#1, sharp corners]
}{%
    \end{tcolorbox}
}

\usepackage{tcolorbox}
\tcbuselibrary{skins,breakable}

\begin{document}

% ============================================================
% PAGE DE GARDE
% ============================================================
\begin{titlepage}
    \centering
    \vspace*{1cm}
    
    % Logo fictif
    \begin{tikzpicture}
        \draw[fill=primaryBlue, rounded corners] (0,0) rectangle (4,1.5);
        \node[white, font=\Large\bfseries] at (2,0.75) {RTSP Stream};
    \end{tikzpicture}
    
    \vspace{2cm}
    
    {\Huge\bfseries\color{primaryBlue} RAPPORT TECHNIQUE\\[0.3cm]}
    {\huge\bfseries SYSTÈME MULTI-RTSP DYNAMIQUE\\[0.5cm]}
    {\Large Serveur C++ (GStreamer) + Interface Python (PyQt)\\[0.3cm]
     Rechargement dynamique}
    
    \vspace{2cm}
    
    \begin{tcolorbox}[colback=lightGray!30, colframe=primaryBlue, width=0.7\textwidth, sharp corners]
        \begin{tabular}{p{3cm}p{8cm}}
            \textbf{Réalisé par} & \textbf{Ayoub Smayen} \\
            \textbf{Fonction} & Développeur C++ \\
            \textbf{Spécialisation} & Streaming vidéo, architectures embarquées \\
            \textbf{Technologies} & GStreamer, C++17, Python, PyQt, OpenCV \\
            \textbf{Version} & 2.0 \\
            \textbf{Date} & Avril 2026 \\
        \end{tabular}
    \end{tcolorbox}
    
    \vfill
    
    \begin{center}
        \rule{0.8\textwidth}{0.5pt}\\
        \textit{Document confidentiel - Tous droits réservés}
    \end{center}
\end{titlepage}

% ============================================================
% TABLE DES MATIERES
% ============================================================
\tableofcontents
\newpage
\listoffigures
\newpage
\listoftables
\newpage

% ============================================================
% CHAPITRE 1 : INTRODUCTION
% ============================================================
\chapter{Introduction}

\section{Contexte et objectifs}

Le protocole \important{RTSP (Real Time Streaming Protocol)} est un standard de l'industrie pour la diffusion de médias en temps réel. Ce projet consiste à développer une solution complète de diffusion et visualisation multi-flux RTSP.

\begin{infobox}{Objectifs du projet}
\begin{itemize}
    \item Diffusion de flux vidéo multi-canal sur réseau local ou Internet
    \item Compatibilité avec les lecteurs standards (VLC, FFplay, OBS)
    \item Configuration dynamique \success{sans interruption du service}
    \item Interface d'administration intuitive
    \item Incrustation de texte personnalisée par canal
\end{itemize}
\end{infobox}

\section{Architecture globale}

Le système se décompose en deux parties indissociables :

\begin{enumerate}
    \item \important{Serveur RTSP multi-canal} (C++17 + GStreamer)
    \begin{itemize}
        \item Diffusion de flux vidéo (tests, fichiers, caméras)
        \item Rechargement à chaud de la configuration
        \item Incrustation texte et horodatage par canal
    \end{itemize}
    
    \item \important{Interface de visualisation dynamique} (Python / PyQt + OpenCV)
    \begin{itemize}
        \item Chargement automatique des caméras depuis JSON
        \item Mise à jour des flux sans redémarrage
        \item Thread par caméra pour garantir la fluidité
    \end{itemize}
\end{enumerate}

% ============================================================
% CHAPITRE 2 : ÉTUDE D'EXISTANT
% ============================================================
\chapter{Étude d'existant}

\section{Analyse comparative}

\begin{table}[H]
    \centering
    \caption{Comparaison des solutions existantes}
    \label{tab:existant}
    \rowcolors{2}{lightGray!30}{white}
    \begin{tabular}{|l|c|c|c|c|c|}
        \hline
        \rowcolor{primaryBlue!80}
        \multicolumn{1}{|c|}{\color{white}\textbf{Solution}} & 
        \color{white}\textbf{RTSP} & 
        \color{white}\textbf{Multi-flux} & 
        \color{white}\textbf{UI} & 
        \color{white}\textbf{Reconfigurable} & 
        \color{white}\textbf{Customisable} \\
        \hline
        VLC Media Player & \success{✅} & ❌ & \success{✅} & ❌ & ❌ \\
        OpenCV seul & \success{✅} & \success{✅} & ❌ & ❌ & \success{✅} \\
        ZoneMinder & \success{✅} & \success{✅} & \success{✅} & ❌ & ❌ \\
        FFmpeg + SDL & \success{✅} & ❌ & ❌ & ❌ & \success{✅} \\
        OBS Studio & \success{✅} & ❌ & \success{✅} & ❌ & ❌ \\
        \hline
        \rowcolor{accentGreen!20}
        \textbf{Notre solution} & \success{✅} & \success{✅} & \success{✅} & \success{✅} & \success{✅} \\
        \hline
    \end{tabular}
\end{table}

\section{Conclusion}

Aucune solution open-source simple ne combine :
\begin{itemize}
    \item Multi-RTSP temps réel
    \item Interface utilisateur moderne
    \item Rechargement dynamique sans coupure
    \item Extensibilité (C++ pour le cœur, Python pour l'UI)
\end{itemize}

➡️ \important{D'où la pertinence de notre architecture hybride.}

% ============================================================
% CHAPITRE 3 : DIAGRAMME DE CAS D'UTILISATION
% ============================================================
\chapter{Diagramme de Cas d'Utilisation}

\section{Diagramme UML}

\begin{figure}[H]
    \centering
    \begin{tikzpicture}[
        actor/.style={rectangle, draw=primaryBlue, fill=primaryBlue!10, minimum width=2cm, minimum height=0.8cm, rounded corners, font=\bfseries},
        usecase/.style={rectangle, draw=secondaryBlue, fill=secondaryBlue!10, minimum width=3.5cm, minimum height=0.7cm, rounded corners, font=\small},
        system/.style={rectangle, draw=primaryBlue, fill=white, minimum width=10cm, minimum height=7cm, rounded corners, dashed},
        include/.style={->, >=stealth, thick, draw=accentOrange, dashed},
        extend/.style={->, >=stealth, thick, draw=accentGreen, dashed}
    ]
        
        % Système
        \draw[system] (-5,-1) rectangle (5,-9);
        \node[font=\large\bfseries, primaryBlue] at (0,-1.3) {Système RTSP Multi-Channel};
        
        % Acteurs
        \node[actor] (admin) at (-7,-2) {Administrateur};
        \node[actor] (viewer) at (-7,-5) {Spectateur};
        \node[actor] (fs) at (7,-3) {Système Fichier};
        
        % Cas d'utilisation
        \node[usecase] (start) at (-2,-2.5) {Démarrer serveur};
        \node[usecase] (add) at (-2,-3.8) {Ajouter canal};
        \node[usecase] (edit) at (-2,-5.1) {Modifier canal};
        \node[usecase] (delete) at (-2,-6.4) {Supprimer canal};
        
        \node[usecase] (watch) at (2,-3.8) {Visualiser flux};
        \node[usecase] (reload) at (2,-5.5) {Rechargement auto};
        
        \node[usecase] (config) at (2,-2.5) {Config JSON};
        
        % Associations Acteur -> UC
        \draw[->, >=stealth, thick] (admin) -- (start);
        \draw[->, >=stealth, thick] (admin) -- (add);
        \draw[->, >=stealth, thick] (admin) -- (edit);
        \draw[->, >=stealth, thick] (admin) -- (delete);
        
        \draw[->, >=stealth, thick] (viewer) -- (watch);
        
        % Associations UC -> Acteur
        \draw[->, >=stealth, thick] (config) -- (fs);
        \draw[->, >=stealth, thick] (reload) -- (fs);
        
        % Inclusions
        \draw[include] (add) -- (config);
        \draw[include] (edit) -- (config);
        \draw[include] (delete) -- (config);
        
        \draw[include] (reload) -- (config);
        
    \end{tikzpicture}
    \caption{Diagramme de cas d'utilisation}
    \label{fig:usecase}
\end{figure}

\section{Description des cas d'utilisation}

\begin{table}[H]
    \centering
    \caption{Cas d'utilisation - Administrateur}
    \label{tab:uc_admin}
    \begin{tabular}{|p{2.5cm}|p{10cm}|}
        \hline
        \rowcolor{primaryBlue!80}
        \multicolumn{2}{|c|}{\color{white}\textbf{UC-01 : Démarrer le serveur RTSP}} \\
        \hline
        \textbf{Acteur principal} & Administrateur \\
        \textbf{Pré-condition} & Exécutable compilé, fichier config valide \\
        \textbf{Post-condition} & Serveur actif, canaux disponibles en RTSP \\
        \textbf{Scénario nominal} & 1. L'admin lance la GUI\\
        & 2. Il clique sur "Démarrer serveur"\\
        & 3. Le serveur C++ se lance\\
        & 4. Les flux deviennent accessibles \\
        \hline
    \end{tabular}
\end{table}

% ============================================================
% CHAPITRE 4 : DIAGRAMME DE SÉQUENCE
% ============================================================
\chapter{Diagramme de Séquence}

\section{Rechargement automatique de la configuration}

\begin{figure}[H]
    \centering
    \begin{tikzpicture}[>=stealth, font=\small]
        
        % Définition des lignes de vie
        \begin{umlseqdiagram}
            \umlactor[class=Administrateur]{admin}
            \umlobject[class=Système Fichier]{fs}
            \umlobject[class=ConfigWatcher (C++)]{watcher}
            \umlobject[class=RTSPServerManager]{server}
            \umlobject[class=GstRTSPMountPoints]{mounts}
            
            \begin{umlcall}[op={modifie rtsp\_config.json}]{admin}{fs}
            \end{umlcall}
            
            \begin{umlcall}[op={événement modification}, type=asynchronous]{fs}{watcher}
            \end{umlcall}
            
            \begin{umlcall}[op={reload\_config()}]{watcher}{server}
                \begin{umlcall}[op={compare(old, new)}]{server}{server}
                \end{umlcall}
            \end{umlcall}
            
            \begin{umlcall}[op={remove\_factory(anciens)}]{server}{mounts}
            \end{umlcall}
            
            \begin{umlcall}[op={add\_factory(nouveaux)}]{server}{mounts}
            \end{umlcall}
            
            \begin{umlcall}[op={return OK}]{mounts}{server}
            \end{umlcall}
            
            \begin{umlcall}[op={log("Config rechargée")}]{server}{watcher}
            \end{umlcall}
            
        \end{umlseqdiagram}
        
    \end{tikzpicture}
    \caption{Diagramme de séquence - Rechargement serveur}
    \label{fig:sequence_server}
\end{figure}

\section{Rechargement côté visualisation PyQt}

\begin{figure}[H]
    \centering
    \begin{tikzpicture}[>=stealth, font=\small]
        
        \begin{umlseqdiagram}
            \umlactor[class=Administrateur]{admin}
            \umlobject[class=watchdog (Python)]{watchdog}
            \umlobject[class=ConfigManager]{config}
            \umlobject[class=MultiRTSPWindow]{window}
            \umlobject[class=CameraStream (Thread)]{thread}
            
            \begin{umlcall}[op={modifie cameras.json}]{admin}{watchdog}
            \end{umlcall}
            
            \begin{umlcall}[op={on\_modified()}, type=asynchronous]{watchdog}{config}
            \end{umlcall}
            
            \begin{umlcall}[op={load()}]{config}{config}
            \end{umlcall}
            
            \begin{umlcall}[op={reload\_config()}]{config}{window}
                \begin{umlcall}[op={compare(old, new)}]{window}{window}
                \end{umlcall}
                
                \begin{umlcall}[op={stop()}, type=asynchronous]{window}{thread}
                \end{umlcall}
                
                \begin{umlcall}[op={start()}, type=asynchronous]{window}{thread}
                \end{umlcall}
                
                \begin{umlcall}[op={update\_UI()}]{window}{window}
                \end{umlcall}
            \end{umlcall}
            
        \end{umlseqdiagram}
        
    \end{tikzpicture}
    \caption{Diagramme de séquence - Rechargement visualisation}
    \label{fig:sequence_viewer}
\end{figure}

% ============================================================
% CHAPITRE 5 : DIAGRAMME DE CLASSES
% ============================================================
\chapter{Diagramme de Classes}

\section{Classes du serveur C++}

\begin{figure}[H]
    \centering
    \begin{tikzpicture}[
        class/.style={rectangle, draw=primaryBlue, fill=primaryBlue!10, minimum width=5cm, align=left, rounded corners, inner sep=5pt},
        attribute/.style={font=\small\itshape},
        method/.style={font=\small},
        relation/.style={-{Triangle[open]}, thick, draw=accentOrange}
    ]
        
        % Classe ChannelConfig
        \node[class] (config) at (0,0) {
            \textbf{<<struct>> ChannelConfig}\\
            \hline
            \attribute{- id : int}\\
            \attribute{- name : string}\\
            \attribute{- mount : string}\\
            \attribute{- source : string}\\
            \attribute{- text\_overlay : string}\\
            \attribute{- enabled : bool}\\
        };
        
        % Classe RTSPServerManager
        \node[class] (manager) at (6,-3) {
            \textbf{RTSPServerManager}\\
            \hline
            \attribute{- server\_ : GstRTSPServer*}\\
            \attribute{- mounts\_ : GstRTSPMountPoints*}\\
            \attribute{- port : int}\\
            \attribute{- host : string}\\
            \hline
            \method{+ init() : bool}\\
            \method{+ add\_channel(cfg) : void}\\
            \method{+ remove\_channel(mount) : void}\\
            \method{+ load\_channels() : void}\\
            \method{+ start() : void}\\
            \method{+ stop() : void}\\
        };
        
        % Classe ConfigWatcher
        \node[class] (watcher) at (-5,-3) {
            \textbf{ConfigWatcher}\\
            \hline
            \attribute{- path : string}\\
            \attribute{- last\_mtime : time\_t}\\
            \attribute{- running : bool}\\
            \hline
            \method{+ start() : void}\\
            \method{+ stop() : void}\\
            \method{- check\_update() : void}\\
        };
        
        % Relations
        \draw[relation] (manager) -- node[above, sloped] {utilise 1..*} (config);
        \draw[relation] (watcher) -- node[above, sloped] {notifie} (manager);
        
    \end{tikzpicture}
    \caption{Diagramme de classes - Serveur C++}
    \label{fig:classes_cpp}
\end{figure}

\section{Classes de l'interface visualisation (PyQt)}

\begin{figure}[H]
    \centering
    \begin{tikzpicture}[
        class/.style={rectangle, draw=secondaryBlue, fill=secondaryBlue!10, minimum width=6cm, align=left, rounded corners, inner sep=5pt},
        qtclass/.style={rectangle, draw=accentGreen, fill=accentGreen!10, minimum width=6cm, align=left, rounded corners, inner sep=5pt},
        relation/.style={-{Triangle[open]}, thick, draw=accentOrange},
        composition/.style={-{Diamond[fill=black]}, thick, draw=primaryBlue}
    ]
        
        % Classe CameraStream
        \node[qtclass] (camera) at (0,0) {
            \textbf{CameraStream (QThread)}\\
            \hline
            \attribute{- url : str}\\
            \attribute{- frame : np.ndarray}\\
            \attribute{- running : bool}\\
            \hline
            \method{+ run() : void}\\
            \method{+ stop() : void}\\
            \method{+ reconnect() : void}\\
        };
        
        % Classe ConfigManager
        \node[class] (config) at (-5,-4) {
            \textbf{ConfigManager}\\
            \hline
            \attribute{- path : str}\\
            \attribute{- data : dict}\\
            \hline
            \method{+ load() : dict}\\
            \method{+ save() : bool}\\
            \method{+ compare(old,new) : bool}\\
        };
        
        % Classe MultiRTSPWindow
        \node[class] (window) at (4,-4) {
            \textbf{MultiRTSPWindow (QMainWindow)}\\
            \hline
            \attribute{- cameras : dict}\\
            \attribute{- labels : dict}\\
            \attribute{- config\_mgr : ConfigManager}\\
            \hline
            \method{+ load\_config() : void}\\
            \method{+ update\_frames() : void}\\
            \method{+ check\_config\_update() : void}\\
            \method{- add\_camera\_widget() : void}\\
            \method{- remove\_camera\_widget() : void}\\
        };
        
        % Classe ConfigFileHandler
        \node[class] (handler) at (-5,-7) {
            \textbf{ConfigFileHandler}\\
            \hline
            \attribute{- path : str}\\
            \attribute{- callback : callable}\\
            \hline
            \method{+ on\_modified(event) : void}\\
        };
        
        % Relations
        \draw[composition] (window) -- node[above, sloped] {1..*} (camera);
        \draw[relation] (window) -- node[above, sloped] {1} (config);
        \draw[relation] (handler) -- node[above, sloped] {notifie} (config);
        
    \end{tikzpicture}
    \caption{Diagramme de classes - Visualisation PyQt}
    \label{fig:classes_python}
\end{figure}

% ============================================================
% CHAPITRE 6 : ARCHITECTURE
% ============================================================
\chapter{Architecture du Système}

\section{Vue d'ensemble}

\begin{figure}[H]
    \centering
    \begin{tikzpicture}[
        block/.style={rectangle, draw=primaryBlue, fill=primaryBlue!10, minimum width=3cm, minimum height=1cm, align=center, rounded corners, drop shadow},
        component/.style={rectangle, draw=secondaryBlue, fill=secondaryBlue!10, minimum width=2.5cm, minimum height=0.8cm, align=center, rounded corners},
        arrow/.style={->, >=stealth, thick, draw=accentOrange},
        datalink/.style={<->, >=stealth, thick, draw=accentGreen}
    ]
        
        % Clients
        \node[block] (vlc) at (0,5) {Client VLC\\\small Spectateur};
        \node[block] (pyqt) at (4,5) {Client PyQt\\\small Visualisateur};
        
        % Serveur
        \node[block] (server) at (2,2.5) {Serveur RTSP C++\\\small GStreamer};
        
        % Configuration
        \node[block] (config) at (2,0) {Fichier JSON\\\small Configuration};
        
        % GUI Admin
        \node[block] (gui) at (-3,0) {GUI Tkinter\\\small Administration};
        
        % Sources
        \node[component] (source1) at (-2,4) {Source Test};
        \node[component] (source2) at (0,4) {Fichier vidéo};
        \node[component] (source3) at (2,4) {Caméra IP};
        \node[component] (source4) at (4,4) {Flux RTSP externe};
        
        % Flèches
        \draw[arrow] (source1) -- (server);
        \draw[arrow] (source2) -- (server);
        \draw[arrow] (source3) -- (server);
        \draw[arrow] (source4) -- (server);
        
        \draw[arrow] (server) -- (vlc);
        \draw[arrow] (server) -- (pyqt);
        
        \draw[datalink] (server) -- (config);
        \draw[datalink] (gui) -- (config);
        
    \end{tikzpicture}
    \caption{Architecture générale du système}
    \label{fig:architecture}
\end{figure}

\section{Pipeline GStreamer}

\begin{figure}[H]
    \centering
    \begin{tikzpicture}[
        element/.style={rectangle, draw=primaryBlue, fill=primaryBlue!15, minimum height=0.8cm, align=center, rounded corners, font=\small},
        arrow/.style={->, >=stealth, thick, draw=accentOrange}
    ]
        
        \node[element] (src) at (0,0) {videotestsrc\\\tiny pattern=0};
        \node[element] (caps) at (2.2,0) {caps filter\\\tiny video/x-raw};
        \node[element] (clock) at (4.4,0) {clockoverlay\\\tiny horodatage};
        \node[element] (text) at (6.6,0) {textoverlay\\\tiny texte personnalisé};
        \node[element] (conv) at (8.8,0) {videoconvert};
        \node[element] (enc) at (11,0) {x264enc\\\tiny H.264};
        \node[element] (pay) at (13.2,0) {rtph264pay\\\tiny RTP};
        
        \draw[arrow] (src) -- (caps);
        \draw[arrow] (caps) -- (clock);
        \draw[arrow] (clock) -- (text);
        \draw[arrow] (text) -- (conv);
        \draw[arrow] (conv) -- (enc);
        \draw[arrow] (enc) -- (pay);
        
    \end{tikzpicture}
    \caption{Pipeline GStreamer par canal}
    \label{fig:pipeline}
\end{figure}

% ============================================================
% CHAPITRE 7 : INSTALLATION
% ============================================================
\chapter{Installation et Déploiement}

\section{Prérequis}

\begin{table}[H]
    \centering
    \caption{Prérequis par plateforme}
    \label{tab:prerequis}
    \begin{tabular}{|l|l|l|l|}
        \hline
        \rowcolor{primaryBlue!80}
        \multicolumn{1}{|c|}{\color{white}\textbf{Composant}} & 
        \color{white}\textbf{Windows} & 
        \color{white}\textbf{Linux} & 
        \color{white}\textbf{macOS} \\
        \hline
        GStreamer & \code{gstreamer-1.0.msi} & \code{apt install} & \code{brew install} \\
        CMake & \code{cmake-3.16+.msi} & \code{apt install} & \code{brew install} \\
        Python & \code{python-3.10+.exe} & \code{apt install} & \code{brew install} \\
        OpenCV & \code{pip install} & \code{pip install} & \code{pip install} \\
        PyQt5 & \code{pip install} & \code{pip install} & \code{pip install} \\
        watchdog & \code{pip install} & \code{pip install} & \code{pip install} \\
        \hline
    \end{tabular}
\end{table}

\section{Compilation}

\subsection{Linux / macOS}
\begin{lstlisting}[language=bash, caption=Compilation sous Linux/macOS]
# Installation des dépendances
sudo apt install libgstreamer1.0-dev libgstreamer-rtsp-server-1.0-dev
sudo apt install cmake build-essential python3 python3-pip

# Compilation
cmake -B build
cmake --build build

# Lancement
./build/rtsp_server config/rtsp_config.json
\end{lstlisting}

\subsection{Windows}
\begin{lstlisting}[language=batch, caption=Compilation sous Windows]
cmake -B build -DGST_ROOT="C:/Program Files/gstreamer/1.0/msvc_x86_64"
cmake --build build --config Release

# Lancement
.\build\Release\rtsp_server.exe config\rtsp_config.json
\end{lstlisting}

% ============================================================
% CHAPITRE 8 : TESTS
% ============================================================
\chapter{Plan de Tests}

\begin{table}[H]
    \centering
    \caption{Plan de tests fonctionnels}
    \label{tab:tests}
    \rowcolors{2}{lightGray!20}{white}
    \begin{adjustbox}{max width=\textwidth}
    \begin{tabular}{|c|l|l|c|}
        \hline
        \rowcolor{primaryBlue!80}
        \multicolumn{1}{|c|}{\color{white}\textbf{ID}} & 
        \color{white}\textbf{Cas de test} & 
        \color{white}\textbf{Résultat attendu} & 
        \color{white}\textbf{Statut} \\
        \hline
        T-01 & Démarrer serveur sans config & Message d'erreur & \success{✅ OK} \\
        T-02 & 3 canaux simultanés & URLs accessibles & \success{✅ OK} \\
        T-03 & Modification du port & Redémarrage sur nouveau port & \success{✅ OK} \\
        T-04 & Ajout canal via GUI & Visible en < 3s & \success{✅ OK} \\
        T-05 & Désactivation canal & Non accessible & \success{✅ OK} \\
        T-06 & Modification overlay texte & Changement instantané & \success{✅ OK} \\
        T-07 & 10 flux simultanés & Stable < 200ms & \success{✅ OK} \\
        T-08 & Ouverture VLC depuis GUI & VLC lance le flux & \success{✅ OK} \\
        T-09 & Copie URL & URL dans presse-papier & \success{✅ OK} \\
        T-10 & Arrêt propre & Clients déconnectés & \success{✅ OK} \\
        T-11 & Rechargement config & Changements < 3s & \success{✅ OK} \\
        T-12 & Chargement cameras.json & Flux affichés & \success{✅ OK} \\
        T-13 & PyQt - Ajout caméra & Nouveau flux apparaît & \success{✅ OK} \\
        T-14 & Reconnexion auto & Reconnexion après perte & \success{✅ OK} \\
        \hline
    \end{tabular}
    \end{adjustbox}
\end{table}

% ============================================================
% CHAPITRE 9 : CONCLUSION
% ============================================================
\chapter{Conclusion}

\section{Bilan}

Ce projet apporte une \important{réponse complète et innovante} au besoin de diffusion et de visualisation multi-flux RTSP avec reconfiguration dynamique.

\section{Points forts}

\begin{itemize}
    \item \success{Séparation des responsabilités} : cœur performant en C++, interface flexible en Python
    \item \success{Rechargement à chaud} des deux côtés (serveur et visualisation)
    \item \success{Compatibilité standards} : VLC, FFplay, OBS
    \item \success{Extensibilité} : facile d'ajouter des fonctionnalités
    \item \success{Multiplateforme} : Windows, Linux, macOS
\end{itemize}

\section{Apport personnel}

\begin{infobox}{Ayoub Smayen – Développeur C++}
En tant que développeur C++, j'ai conçu et implémenté :
\begin{itemize}
    \item Le serveur RTSP basé sur \code{gst-rtsp-server}
    \item Le watcher de fichier config avec rechargement sans interruption
    \item Le pipeline GStreamer optimisé pour la faible latence
    \item L'interface Python/PyQt pour la visualisation dynamique
    \item L'interface Tkinter pour l'administration
\end{itemize}
\end{infobox}

\section{Évolutions prévues}

\begin{table}[H]
    \centering
    \begin{tabular}{|l|l|}
        \hline
        \rowcolor{primaryBlue!80}
        \color{white}\textbf{Fonctionnalité} & \color{white}\textbf{Statut} \\
        \hline
        Support caméras IP natives (rtspsrc) & \warning{�� Prévu} \\
        Authentification RTSP & \warning{�� Prévu} \\
        Enregistrement des flux & \warning{�� Prévu} \\
        Interface web (Flask + WebSockets) & \warning{�� Prévu} \\
        Détection de mouvement (OpenCV) & \warning{�� Prévu} \\
        Notification en cas de perte de flux & \warning{�� Prévu} \\
        \hline
    \end{tabular}
\end{table}

% ============================================================
% ANNEXES
% ============================================================
\chapter*{Annexes}
\addcontentsline{toc}{chapter}{Annexes}

\section*{Annexe A : Extrait de code – Serveur C++}
\addcontentsline{toc}{section}{Annexe A : Extrait de code – Serveur C++}

\begin{lstlisting}[language=c++, caption=RTSPServerManager.cpp]
void RTSPServerManager::add_channel(const ChannelConfig& cfg) {
    if (!cfg.enabled) return;
    
    GstRTSPMediaFactory* factory = gst_rtsp_media_factory_new();
    
    std::string pipeline = cfg.source;
    if (cfg.source == "videotestsrc") {
        pipeline += " pattern=" + std::to_string(cfg.pattern);
    }
    pipeline += " ! video/x-raw ! clockoverlay ! textoverlay text=\"" + 
                cfg.text_overlay + "\" ! videoconvert ! x264enc ! rtph264pay name=pay0 pt=96";
    
    gst_rtsp_media_factory_set_launch(factory, pipeline.c_str());
    gst_rtsp_media_factory_set_shared(factory, TRUE);
    
    gst_rtsp_mount_points_add_factory(mounts_, cfg.mount.c_str(), factory);
    g_object_unref(factory);
    
    g_print("✅ Canal ajoute : %s -> %s\n", cfg.name.c_str(), cfg.mount.c_str());
}
\end{lstlisting}

\section*{Annexe B : Extrait de code – Visualisation PyQt}
\addcontentsline{toc}{section}{Annexe B : Extrait de code – Visualisation PyQt}

\begin{lstlisting}[language=python, caption=camera_stream.py]
import cv2
from PyQt5.QtCore import QThread, pyqtSignal, QMutex

class CameraStream(QThread):
    frame_updated = pyqtSignal(object)
    
    def __init__(self, url: str, parent=None):
        super().__init__(parent)
        self.url = url
        self.running = True
        self.mutex = QMutex()
        
    def run(self):
        cap = cv2.VideoCapture(self.url)
        while self.running:
            ret, frame = cap.read()
            if ret:
                self.frame_updated.emit(frame)
            else:
                cap.release()
                self.msleep(1000)
                cap = cv2.VideoCapture(self.url)
        cap.release()
    
    def stop(self):
        self.mutex.lock()
        self.running = False
        self.mutex.unlock()
\end{lstlisting}

\section*{Annexe C : Fichier de configuration JSON}
\addcontentsline{toc}{section}{Annexe C : Fichier de configuration JSON}

\begin{lstlisting}[language=json, caption=rtsp_config.json]
{
  "server": {
    "host": "0.0.0.0",
    "port": 8554
  },
  "channels": [
    {
      "id": 1,
      "name": "Camera_Principale",
      "mount": "/main",
      "source": "videotestsrc",
      "pattern": 18,
      "text_overlay": "Live - Camera Principale",
      "enabled": true
    },
    {
      "id": 2,
      "name": "Camera_Secondaire",
      "mount": "/secondary",
      "source": "filesrc location=/home/video/test.mp4 ! decodebin",
      "pattern": 0,
      "text_overlay": "Replay - Camera Secondaire",
      "enabled": true
    }
  ]
}
\end{lstlisting}

% ============================================================
% FIN
% ============================================================
\end{document}